% English translation of checked_pages/esq_119_005.tex; source: cote 119, site 5, PDF page 6, typescript page 3.
% AI-translated by GPT-6 Astra (OpenAI, ChatGPT Pro mode).
\scanpage{119/5}
\mspage{3}

1970, my attention was systematically directed toward objects
of the greatest generality, in order to draw out an adequate overall language
for the world of algebraic geometry, and I lingered
over algebraic curves only to the strict extent that this proved
indispensable (particularly in étale cohomology) for developing
``all-purpose'' techniques and statements valid in every dimension
and in every place (I mean, over all base schemes, or even
all ringed base topoi ...), here I am thus brought back, through
objects so simple that a child can come to know them through play,
to the beginnings and origins of algebraic geometry, familiar to Riemann
and his followers !

Since about 1975, then, the geometry of (real)
surfaces, and from 1977 onward the connections between questions in
the geometry of surfaces and the algebraic geometry of algebraic curves
defined over fields such as $\mathbb{C}$, $\mathbb{R}$ or extensions of finite
type of $\mathbb{Q}$, have been my principal source of inspiration, as well
as my constant guiding thread. It was with surprise and wonder
that over the years I discovered (or rather, probably,
rediscovered) the prodigious, truly inexhaustible richness, the
unsuspected depth of this seemingly so innocuous theme. I believe
I sense in it a particularly vital point, a privileged point of convergence
of the main currents of mathematical ideas, as also
of the main structures and the visions of things they express,
from the most specific (such as the rings $\mathbb{Z}$, $\mathbb{Q}$, $\bar{\mathbb{Q}}$, $\mathbb{R}$, $\mathbb{C}$,
or the group $S\ell(2)$ over one of these rings, or general reductive
algebraic groups) to the most ``abstract'', such as algebraic,
complex analytic or real analytic ``multiplicities''.\tnote{TN-005-01}{``Multiplicity'' consistently renders the author's ``multiplicité''. His parenthesis here explains its role in retaining the automorphisms of the objects under study.}
(These enter naturally when one sets out to study
``moduli varieties'' systematically for the geometric objects
under consideration, seeking to go beyond the notoriously insufficient
viewpoint of ``coarse moduli'', which amounts, most unfortunately,
to ``killing'' the automorphism groups of these objects .) Among these
moduli multiplicities, it is those of Mumford-Deligne for
``stable'' algebraic curves of genus $g$, with $\nu$ marked points, which

\margin{5}

